\documentclass[11pt,a4paper]{article}
\usepackage[utf8]{inputenc}
\usepackage[T1]{fontenc}
\usepackage[portuguese]{babel}
\usepackage{xcolor}
\usepackage{hyperref}
\usepackage{geometry}
\usepackage{tcolorbox}

\geometry{margin=2.5cm}

\definecolor{primary}{RGB}{39, 174, 96}
\definecolor{secondary}{RGB}{44, 62, 80}
\definecolor{codebg}{RGB}{240, 240, 240}

\hypersetup{colorlinks=true, linkcolor=primary, urlcolor=primary}

\title{\color{secondary}\Huge\textbf{Solução: Exercício 4.5 - O projeto, passo 22}}
\author{\color{primary}DevOps com Kubernetes -- 2026}
\date{}

\begin{document}
\maketitle

\section*{\color{primary}Guia de Solução}
\begin{tcolorbox}[colback=green!5!white,colframe=green!60!black,title=Resolução Passo a Passo]
### Passo a Passo

1. \textbf{Modificação no Banco de Dados}: Adicione uma coluna/campo booleano \texttt{done} (concluído) na sua tabela ou esquema de tarefas (Todos), com valor padrão \texttt{false}.
2. \textbf{Atualizar o Backend}: Crie o endpoint \texttt{PUT /todos/:id} no backend. O controlador dessa rota deve buscar a tarefa correspondente no banco de dados e atualizar o valor de \texttt{done} para \texttt{true}.
3. \textbf{Atualizar o Frontend}: Na interface de usuário, adicione um botão 'Concluído' ou um \textit{checkbox} ao lado de cada tarefa listada. Ao ser clicado, o botão envia a requisição HTTP PUT para a nova rota da API.
4. \textbf{Implantação}: Atualize as imagens Docker, altere as tags no \textit{Deployment} e aplique no Kubernetes. Valide a funcionalidade criando uma tarefa e depois marcando-a como concluída.
\end{tcolorbox}

\end{document}
